% !TEX root =../tfmETSI.tex
\chapter{Conclusions and Future Work}\label{chp:conclusion}

\section{Conclusions}
\label{sec:concl-conclusions}
\begin{itemize}
    \item Baseline result (RQ1): 63.6\% / 62.7\% macro-F1, worst class lung adenocarcinoma.
    \item Pathology pretraining (RQ2): PLIP $+7.3$ macro-F1, mainly benign vs malignant.
    \item Feature adaptation (RQ3): partial unfreezing (last 30 of 178 layers) $+27.1$ macro-F1, broad-based across classes.
    \item CLIP-paradigm value (RQ4): CLIP and plain CE tie at the frozen-features ceiling ($\sim 0.63$); CLIP wins $+2.6$ macro-F1 over CE under feature adaptation (unfreeze30), broad-based across all classes. The paradigm's primary value remains in capabilities (open-vocabulary, retrieval), not raw accuracy.
    \item Stain normalization (RQ5): the gain is strongly method-specific. Macenko yields $+11.7$ macro-F1 ($0.627 \to 0.739$) at the frozen-backbone budget, larger than the PLIP zero-shot gain; Vahadane yields $+3.7$ and Reinhard $+2.5$ on the identical pipeline. Composing the best stain method (Macenko) with the best feature-adaptation regime (unfreeze30) reaches \textbf{0.918 macro-F1}, the highest measured in this study --- $+0.020$ over unfreeze30 alone, so the levers compose with diminishing returns rather than additively.
    \item Biomedical text encoder (RQ6): \textit{TBD}.
    \item OOD transfer (RQ7): the picture is two-level. The best in-distribution model (Macenko + Unfreeze30, $0.918$) drops to $0.000$ macro-F1 when the same Macenko transform is applied at OOD time --- NCT is already colour-normalised against its own reference and the double-shift is fatal. With the OOD-time Macenko stripped, the same checkpoint reaches $0.201$. The raw-trained Unfreeze30 is the best LC25000-trained transfer at $0.331$; \textbf{PLIP zero-shot beats every LC25000-trained variant at $0.661$ macro-F1}. The cross-distribution UMAP (Figure~\ref{fig:ood-cross-umap}) reveals that LC25000 fine-tuning \emph{does} transfer organ-level structure (NCT colon patches land in the LC25000 colon region of embedding space, not in lung), so the predictions stay within the colon classes; what does not transfer is the within-colon benign-vs-malignant discrimination, which depends on cellular cues that are more sensitive to lab/scanner/label-granularity differences. Only pathology pretraining (PLIP) produces representations that transfer at both organ \emph{and} tissue-state level.
    \item Cross-cutting observations: saturation of frozen transfer, role of dataset artifacts.
\end{itemize}


\section{Limitations}
\label{sec:concl-limitations}
\begin{itemize}
    \item Single-institution dataset, no patient-level grouping, no clinical validation.
    \item Small batch + few classes $\to$ InfoNCE false negatives, documented not mitigated.
    \item Frozen-backbone baseline tests transfer, not full fine-tuning.
    \item No human-expert evaluation.
    \item No head-to-head against CONCH / BiomedCLIP / MUSK (compute scope).
\end{itemize}


\section{Future Work}
\label{sec:concl-future}

Near-term:
\begin{itemize}
    \item Additional external datasets beyond colon-Kaggle.
    \item SupCon or class-balanced sampling for false-negative regime.
    \item Patient-level grouping if metadata recoverable.
\end{itemize}

Not run in this project (low expected $\Delta$ under frozen text encoders + short prompts):
\begin{itemize}
    \item Inference-time prompt ensembling.
    \item Dynamic prompt training.
\end{itemize}

Medium-term:
\begin{itemize}
    \item Larger effective batch (grad accumulation or bigger GPU).
    \item Prompt-tuning methods (CoOp, CoCoOp).
\end{itemize}

Longer-term:
\begin{itemize}
    \item Architecture-controlled comparison against CONCH / BiomedCLIP / MUSK / PathChat.
    \item Multi-institution training and evaluation.
    \item Returning to segmentation with pathology FM backbones.
    \item CLIP-with-short-labels applied to other low-label scientific domains.
\end{itemize}
